% Options for packages loaded elsewhere
% Options for packages loaded elsewhere
\PassOptionsToPackage{unicode}{hyperref}
\PassOptionsToPackage{hyphens}{url}
\PassOptionsToPackage{dvipsnames,svgnames,x11names}{xcolor}
%
\documentclass[
  letterpaper,
  DIV=11,
  numbers=noendperiod]{scrartcl}
\usepackage{xcolor}
\usepackage{amsmath,amssymb}
\setcounter{secnumdepth}{3}
\usepackage{iftex}
\ifPDFTeX
  \usepackage[T1]{fontenc}
  \usepackage[utf8]{inputenc}
  \usepackage{textcomp} % provide euro and other symbols
\else % if luatex or xetex
  \usepackage{unicode-math} % this also loads fontspec
  \defaultfontfeatures{Scale=MatchLowercase}
  \defaultfontfeatures[\rmfamily]{Ligatures=TeX,Scale=1}
\fi
\usepackage{lmodern}
\ifPDFTeX\else
  % xetex/luatex font selection
\fi
% Use upquote if available, for straight quotes in verbatim environments
\IfFileExists{upquote.sty}{\usepackage{upquote}}{}
\IfFileExists{microtype.sty}{% use microtype if available
  \usepackage[]{microtype}
  \UseMicrotypeSet[protrusion]{basicmath} % disable protrusion for tt fonts
}{}
\makeatletter
\@ifundefined{KOMAClassName}{% if non-KOMA class
  \IfFileExists{parskip.sty}{%
    \usepackage{parskip}
  }{% else
    \setlength{\parindent}{0pt}
    \setlength{\parskip}{6pt plus 2pt minus 1pt}}
}{% if KOMA class
  \KOMAoptions{parskip=half}}
\makeatother
% Make \paragraph and \subparagraph free-standing
\makeatletter
\ifx\paragraph\undefined\else
  \let\oldparagraph\paragraph
  \renewcommand{\paragraph}{
    \@ifstar
      \xxxParagraphStar
      \xxxParagraphNoStar
  }
  \newcommand{\xxxParagraphStar}[1]{\oldparagraph*{#1}\mbox{}}
  \newcommand{\xxxParagraphNoStar}[1]{\oldparagraph{#1}\mbox{}}
\fi
\ifx\subparagraph\undefined\else
  \let\oldsubparagraph\subparagraph
  \renewcommand{\subparagraph}{
    \@ifstar
      \xxxSubParagraphStar
      \xxxSubParagraphNoStar
  }
  \newcommand{\xxxSubParagraphStar}[1]{\oldsubparagraph*{#1}\mbox{}}
  \newcommand{\xxxSubParagraphNoStar}[1]{\oldsubparagraph{#1}\mbox{}}
\fi
\makeatother


\usepackage{longtable,booktabs,array}
\newcounter{none} % for unnumbered tables
\usepackage{calc} % for calculating minipage widths
% Correct order of tables after \paragraph or \subparagraph
\usepackage{etoolbox}
\makeatletter
\patchcmd\longtable{\par}{\if@noskipsec\mbox{}\fi\par}{}{}
\makeatother
% Allow footnotes in longtable head/foot
\IfFileExists{footnotehyper.sty}{\usepackage{footnotehyper}}{\usepackage{footnote}}
\makesavenoteenv{longtable}
\usepackage{graphicx}
\makeatletter
\newsavebox\pandoc@box
\newcommand*\pandocbounded[1]{% scales image to fit in text height/width
  \sbox\pandoc@box{#1}%
  \Gscale@div\@tempa{\textheight}{\dimexpr\ht\pandoc@box+\dp\pandoc@box\relax}%
  \Gscale@div\@tempb{\linewidth}{\wd\pandoc@box}%
  \ifdim\@tempb\p@<\@tempa\p@\let\@tempa\@tempb\fi% select the smaller of both
  \ifdim\@tempa\p@<\p@\scalebox{\@tempa}{\usebox\pandoc@box}%
  \else\usebox{\pandoc@box}%
  \fi%
}
% Set default figure placement to htbp
\def\fps@figure{htbp}
\makeatother


% definitions for citeproc citations
\NewDocumentCommand\citeproctext{}{}
\NewDocumentCommand\citeproc{mm}{%
  \begingroup\def\citeproctext{#2}\cite{#1}\endgroup}
\makeatletter
 % allow citations to break across lines
 \let\@cite@ofmt\@firstofone
 % avoid brackets around text for \cite:
 \def\@biblabel#1{}
 \def\@cite#1#2{{#1\if@tempswa , #2\fi}}
\makeatother
\newlength{\cslhangindent}
\setlength{\cslhangindent}{1.5em}
\newlength{\csllabelwidth}
\setlength{\csllabelwidth}{3em}
\newenvironment{CSLReferences}[2] % #1 hanging-indent, #2 entry-spacing
 {\begin{list}{}{%
  \setlength{\itemindent}{0pt}
  \setlength{\leftmargin}{0pt}
  \setlength{\parsep}{0pt}
  % turn on hanging indent if param 1 is 1
  \ifodd #1
   \setlength{\leftmargin}{\cslhangindent}
   \setlength{\itemindent}{-1\cslhangindent}
  \fi
  % set entry spacing
  \setlength{\itemsep}{#2\baselineskip}}}
 {\end{list}}
\usepackage{calc}
\newcommand{\CSLBlock}[1]{\hfill\break\parbox[t]{\linewidth}{\strut\ignorespaces#1\strut}}
\newcommand{\CSLLeftMargin}[1]{\parbox[t]{\csllabelwidth}{\strut#1\strut}}
\newcommand{\CSLRightInline}[1]{\parbox[t]{\linewidth - \csllabelwidth}{\strut#1\strut}}
\newcommand{\CSLIndent}[1]{\hspace{\cslhangindent}#1}



\setlength{\emergencystretch}{3em} % prevent overfull lines

\providecommand{\tightlist}{%
  \setlength{\itemsep}{0pt}\setlength{\parskip}{0pt}}



 


\KOMAoption{captions}{tableheading}
\makeatletter
\@ifpackageloaded{caption}{}{\usepackage{caption}}
\AtBeginDocument{%
\ifdefined\contentsname
  \renewcommand*\contentsname{Table of contents}
\else
  \newcommand\contentsname{Table of contents}
\fi
\ifdefined\listfigurename
  \renewcommand*\listfigurename{List of Figures}
\else
  \newcommand\listfigurename{List of Figures}
\fi
\ifdefined\listtablename
  \renewcommand*\listtablename{List of Tables}
\else
  \newcommand\listtablename{List of Tables}
\fi
\ifdefined\figurename
  \renewcommand*\figurename{Figure}
\else
  \newcommand\figurename{Figure}
\fi
\ifdefined\tablename
  \renewcommand*\tablename{Table}
\else
  \newcommand\tablename{Table}
\fi
}
\@ifpackageloaded{float}{}{\usepackage{float}}
\floatstyle{ruled}
\@ifundefined{c@chapter}{\newfloat{codelisting}{h}{lop}}{\newfloat{codelisting}{h}{lop}[chapter]}
\floatname{codelisting}{Listing}
\newcommand*\listoflistings{\listof{codelisting}{List of Listings}}
\makeatother
\makeatletter
\makeatother
\makeatletter
\@ifpackageloaded{caption}{}{\usepackage{caption}}
\@ifpackageloaded{subcaption}{}{\usepackage{subcaption}}
\makeatother
\usepackage{bookmark}
\IfFileExists{xurl.sty}{\usepackage{xurl}}{} % add URL line breaks if available
\urlstyle{same}
\hypersetup{
  pdftitle={Infinite Worlds and the Two Envelope Paradox},
  pdfauthor={Brian Weatherson},
  colorlinks=true,
  linkcolor={blue},
  filecolor={Maroon},
  citecolor={Blue},
  urlcolor={Blue},
  pdfcreator={LaTeX via pandoc}}


\title{Infinite Worlds and the Two Envelope Paradox}
\author{Brian Weatherson}
\date{2025-10-14}
\begin{document}
\maketitle
\begin{abstract}
There have been a lot of discussions recently about how versions of the
St Petersburg paradox affect theories about welfare in infinite worlds.
This note is about how the two envelope paradox affects things. The
tentative suggestion is that it gives us a reason to give up on an
otherwise very plausible looking principle of state-level strict
dominance.
\end{abstract}


There have been a flurry of interesting recent papers on welfare
aggregation principles in infinite worlds. A small sample includes:

\begin{itemize}
\tightlist
\item
  Hayden Wilkinson, \emph{Infinite Aggregation and Risk} (Wilkinson
  2023)
\item
  Frank Hong and Jeffrey Sanford Russell, \emph{Paradoxes of Infinite
  Aggregation} (Hong and Russell 2025)
\item
  Jake Nebel, \emph{Infinite Ethics and the Limits of Impartiality}
  (Nebel 2025)
\item
  Jeremy Goodman, \emph{Permutation-Invariant Social Welfare Orders Are
  Anonymous} (Goodman 2025)
\item
  Jeremy Goodman and Harvey Lederman, \emph{Maximal Social Welfare
  Relations on Infinite Populations Satisfying Permutation Invariance}
  (Goodman and Lederman 2024)
\end{itemize}

I'm going to be drawing on all of these here, but especially the last
one.\footnote{And it should be noted that a lot of these works are
  downstream of Amanda Askell's PhD Thesis \emph{Pareto Principles in
  Infinite Ethics} (Askell 2018). The positive point of this note is to
  provide another argument for using what Goodman and Lederman call the
  \emph{Sum Preorder} for comparing infinite worlds. But I'll leave most
  of that argument implicit; hopefully I'll return to it in a later
  post.}

Hong and Russell (2025) use variants of the St Petersburg paradox to
show that the following four principles are inconsistent. Notably, they
do so even given the extra assumption that welfare levels for
individuals can only take one of two possible values.\footnote{It's much
  easier to derive contradictions if we allow that individuals can have
  unbounded utility.} I'll adopt that assumption too.\footnote{The
  following principles are mostly direct quotes from Hong and Russell's
  paper, but I've expanded the notation for \textbf{Stochastic
  Compensation} for clarity.}

\begin{description}
\tightlist
\item[(Weak) Ex Ante Pareto]
If \emph{X} is at least as good as \emph{Y} for every individual, then
\emph{X} is at least as good as \emph{Y} overall.
\item[Stochastic Compensation]
For any good \emph{x} for an individual \emph{i}, and any events
\emph{E} and \emph{F} such that Pr(\emph{E}) ≤ Pr(\emph{F}), there is
some good \emph{y} for \emph{i} such that the gamble ⟨\emph{y} if
\emph{F}, 0 otherwise⟩ is at least as good for \emph{i} than the gamble
⟨\emph{x} if \emph{E}, 0 otherwise⟩.
\item[Interpersonal Compensation]
For any allocation \emph{x} for a finite set of individuals \emph{I},
there is an allocation \emph{y} for some finite set of individuals
\emph{J} disjoint from \emph{y} such that \emph{y} is better overall
than \emph{x}.
\item[(Strict) Statewise Dominance]
If the outcome that results from a prospect \emph{X} is strictly better
than the outcome that results from a prospect \emph{Y} in every possible
state, then \emph{X} is strictly better than \emph{Y}.
\end{description}

The first point I'll make is that slight generalisations of the last two
on their own suffice for a contradiction. That makes me rather sceptical
that the first two principles are the problems. Here are the two
principles that I need.

\begin{description}
\tightlist
\item[Risk Spreading]
For some integer \emph{m}, and some probability \emph{p} \textless{} ½,
the following holds. It is better that \emph{m\textsuperscript{n}}
people get 1 for certain, and
\emph{m}\textsuperscript{\emph{n}+2}~‑~\emph{m\textsuperscript{n}}
people have chance \emph{p} of 1 (and chance 1‑\emph{p} of 0), than that
\emph{m}\textsuperscript{\emph{n}+1} people get 1 and everyone else gets
0.
\item[(Strict) Partition Dominance]
If a prospect \emph{X} is strictly better than a prospect \emph{Y}
conditional on each member of a countable partition, then \emph{X} is
strictly better than \emph{Y}.
\end{description}

I take \textbf{Risk Spreading} to be a probabilistic form of
\textbf{Interpersonal Compensation}. It's better that a much larger
group of people have a good chance (nearly ½) of the good, than that a
much smaller group have the good for sure, and everyone misses out. This
would hold given any plausible kind of Additivity principle, and it's if
anything a fairer distribution of goods.

And I take \textbf{(Strict) Partition Dominance} to be motivated by the
same things that motivate \textbf{(Strict) Statewise Dominance}. When we
talk about states in this bit of philosophy/economics, that's always
shorthand. We never really know in real life what the full outcome of
any action will be. We can save the child from the pond, but they might
grow up to be a mass murderer. The typical, and I think correct,
practice is to simply bracket off those considerations. The relevant
`states' are that the child lives, or the child drowns. In reality, each
of those states is another lottery; it's just that the first is a much
much better one. Any principle that reasons from properties of states to
properties of prospects is really reasoning from properties of more
fine-grained prospects to properties of less fine-grained prospects.
From that perspective, \textbf{(Strict) Partition Dominance} is really
just a reformulation of \textbf{(Strict) Statewise Dominance}.

Given this, that these are inconsistent is quite simple, and just
follows the familiar pattern from the literature on the Two Envelope
Paradox.\footnote{For a primer on that, see Jackson, Menzies, and Oppy
  (1994).} There are two coins, a \textcolor{green}{green} coin and a
fair \textcolor{brown}{brown} coin. Unlike the \textcolor{brown}{brown}
coin, the \textcolor{green}{green} coin is not fair;
\textcolor{green}{it} has chance ε of landing \textcolor{green}{heads}.
The \textcolor{green}{green} coin will be flipped repeatedly until
\textcolor{green}{it} lands \textcolor{green}{heads}. Let \emph{n} be
the number of times \textcolor{green}{it} lands
\textcolor{green}{tails}; unless \textcolor{green}{it} never lands
\textcolor{green}{tails}, in which case \emph{n}~=~0. Then the
\textcolor{brown}{brown} coin is flipped once. Assume (as is normal in
these problems) that we have an infinite population,
\emph{p}\textsubscript{1}, \emph{p}\textsubscript{2}, \ldots, with the
numbers being arbitrary. And then assume we have two prospects, which
we'll call \textcolor{blue}{Blue} and \textcolor{red}{Red}. (Assume
\emph{m} is the same large number that witnesses \textbf{Risk
Spreading}, and that \emph{m}~\textgreater~2.)

\begin{description}
\tightlist
\item[\textcolor{blue}{Blue}]
If the \textcolor{brown}{brown} coin lands \textcolor{brown}{heads}, the
first \emph{m\textsuperscript{n}} people get 1, everyone else gets 0.

If the \textcolor{brown}{brown} coin lands \textcolor{brown}{tails}, the
first \emph{m}\textsuperscript{\emph{n}+1} people get 1, everyone else
gets 0.
\item[\textcolor{red}{Red}]
If the \textcolor{brown}{brown} coin lands \textcolor{brown}{heads}, the
first \emph{m}\textsuperscript{\emph{n}+1} people get 1, everyone else
gets 0.

If the \textcolor{brown}{brown} coin lands \textcolor{brown}{tails}, the
first \emph{m\textsuperscript{n}} people get 1, everyone else gets 0.
\end{description}

Let \emph{B} be a random variable equalling the number of people who
will get 1 if \textcolor{blue}{Blue} is chosen, and \emph{R} be a random
variable equalling the number of people who will get 1 if
\textcolor{red}{Red} is chosen. Just look at the possible values for
\emph{B}; what happens with \emph{R} will be entirely parallel. If
\emph{B}~=~1, then \emph{R}~=~\emph{m}, so \textcolor{red}{Red} is
clearly better. If \emph{B}~=~\emph{m}\textsuperscript{\emph{n}+1} for
\emph{n}~⩾~0, then \emph{R}~=~\emph{m\textsuperscript{n}} with
probability 1/(2‑ε), and \emph{R}~=~\emph{m}\textsuperscript{\emph{n}+2}
with probability (1‑ε)/(2‑ε). Provided ε is small enough, that will be
greater than the value \emph{p} in \textbf{Risk Spreading}, so again
\textcolor{red}{Red} will be better.

So conditional on any possible value of \emph{B}, \textcolor{red}{Red}
is better than \textcolor{blue}{Blue}. So by \textbf{(Strict) Statewise
Dominance}, \textcolor{red}{Red} is better than \textcolor{blue}{Blue}.
But an exactly parallel argument shows that conditional on any possible
value of \emph{R}, \textcolor{blue}{Blue} is better than
\textcolor{red}{Red}. So by \textbf{(Strict) Statewise Dominance},
\textcolor{blue}{Blue} is better than \textcolor{red}{Red}.
Contradiction.

We could probably tighten this up, but I think this shows that fairly
weak dominance principles, plus fairly weak compensation principles, are
inconsistent in infinite worlds. This makes me think that the guilty
principle in Hong and Russell's tetralemma will be one of
\textbf{(Strict) Statewise Dominance} and \textbf{Interpersonal
Compensation}. I'm inclined to give up \textbf{(Strict) Statewise
Dominance}, though it's possible that the right lesson to draw from the
two envelope cases is that there is an argument in the style of Nebel
(2025) against impartiality which does not rely on completeness.

\protect\phantomsection\label{refs}
\begin{CSLReferences}{1}{0}
\bibitem[\citeproctext]{ref-Askell2018}
Askell, Amanda. 2018. {``Pareto Principles in Infinite Ethics.''} PhD
thesis, New York University.

\bibitem[\citeproctext]{ref-Goodman2025}
Goodman, Jeremy. 2025. {``Permutation-Invariant Social Welfare Orders
Are Anonymous.''} \emph{Journal of Mathematical Economics} 120.

\bibitem[\citeproctext]{ref-GoodmanLederemanArXiV}
Goodman, Jeremy, and Harvey Lederman. 2024. {``Maximal Social Welfare
Relations on Infinite Populations Satisfying Permutation Invariance.''}
\url{https://arxiv.org/abs/2408.05851v2}. arXiv preprint.

\bibitem[\citeproctext]{ref-HongRussell2025}
Hong, Frank, and Jeffrey Sanford Russell. 2025. {``Paradoxes of Infinite
Aggregation.''} \emph{Noûs} 59 (3): 809--27. doi:
\href{https://doi.org/10.1111/nous.12535}{10.1111/nous.12535}.

\bibitem[\citeproctext]{ref-TwoEnvelope}
Jackson, Frank, Peter Menzies, and Graham Oppy. 1994. {``The Two
Envelope {`Paradox'}.''} \emph{Analysis} 54 (1): 43--45. doi:
\href{https://doi.org/10.1093/analys/54.1.43}{10.1093/analys/54.1.43}.

\bibitem[\citeproctext]{ref-Nebel2025}
Nebel, Jacob M. 2025. {``Infinite Ethics and the Limits of
Impartiality.''} No{û}s. 2025. doi:
\href{https://doi.org/10.1111/nous.70010}{10.1111/nous.70010}.

\bibitem[\citeproctext]{ref-Wilkinson2023}
Wilkinson, Hayden. 2023. {``Infinite Aggregation and Risk.''}
\emph{Australasian Journal of Philosophy} 101 (2): 340--59. doi:
\href{https://doi.org/10.1080/00048402.2021.2013265}{10.1080/00048402.2021.2013265}.

\end{CSLReferences}




\end{document}
